% Nouveau EGA V (cote 25): archive cote 25, site page 102 (PDF page 103).
% typescript. Draft compared: none.
% AI-transcribed by Claude Opus 5.5 (claude-opus-5-5, adaptive thinking, effort high), Message Batches API, batch ega25b (msgbatch_01A7n5YuXubUPNK1zsepi8RG), 2026-11-13; not yet checked by eye.
\scanpage{25/102}
\mspage{5}
% NOTE: the whole page is crossed by three long diagonal pen strokes running from upper right to lower left (top of the page down to "rang" in the line "i.e..E partout de rang"; from "(A ce pro-" down to "1. 2." in "la remarque 1. 2."; from "projectif" in "Soit P' = ..." down to "définis" in the last line). The text is transcribed as live text under them.
de Picard, déduit de la s ection de ce dernier définie par $\underline{O}_P(1)$, est un
\uline{isomorphisme}.

Soit P un fibré projectif sur un corps k. Un Module inversible
$\underline{L}$ sur P est dit \uline{de degré} n si $\underline{L}$ est isomorphe à $\underline{O}_P(n)$; lorsque
dim P$\geq$1, cela détermine n en \xout{t}termes de $\underline{L}$ , mais si dim P $\leq$ \unsure{0|D}
(i.e. p est vide, ou réduit à un point) alors $\underline{L}$ est degré n pour tout n.
Lorsque P est un fibré projectif sur une base quelconque, un Module
inversible $\underline{L}$ sur P est dit de degré n s'il est de degré n sur chaque fibr
% NOTE: the end of the preceding line is cut off at the right edge of the scan
Cela détermine n, sauf si toutes les fibres de P/S sont de dim$\leq$0, alors
$\underline{L}$ est de degré n pour tout n. Dire que $\underline{L}$ est de degré n signifie aussi,
en vertu de 1.1. et 1.2. , que \inserted[typed]{la classe} \xout{\unsure{xximage}} de $\underline{L}$ dans Pic(P) est dans
l'image de \inserted{$\mathrm{Pic}\,S\times\{n\}$ \unsure{par|pr}} l'homomorphisme (x) Pic(S)x$\underline{Z}$ $\longrightarrow$Pic(P) envisagé plus haut,
i.e. que $\underline{L}$ est isomorphe à un Module de la forme $f^*(\underline{M})(n)$ , où $\underline{M}$ est
inversible sur S. D'ailleurs, lorsque les fibres de P sont non vides
i.e..$\underline{E}$ partout de rang $\geq$1, alors $\underline{M}$ est déterminé à isomorphisme unique
\margin{\textbar}près en termes de $\underline{L}$, comme il résulte encore de 1.1. et 1.2. (A ce pro-%
pos, je m'aperçois qu'il y avait lieu d'énoncer 1.1. ainsi, sans hypothèse
sur le rang de $\underline{E}$ : tout $\underline{L}$ inversible sur S peut s e mettre sous la forme
indiquée dans l'énoncé; lorsque les fibres de P sont non vides i.e. le
rang de $\underline{E}$ partout $\neq$ 0 , alors $\underline{M}$ est déterminé à isom unique près par
$\underline{L}$ ; lorsque les fibres de P sont toutes de dim $\geq$1, \inserted{i.e. le rang de $\underline{E}$ partout $\geq 2$,} alors la partition
de S est également déterminée de façon unique par le choix de $\underline{L}$. De cette
façon, la remarque 1. 2. est éliminée et passe dans la démonstration).

Soit P' = P($\underline{E}$') un deuxième fibré projectif, alors la détermina-%
tion de Pic(P) permet en prihcipe de déterminer les S-morphismes
g:P$\to$P', puisque ceux-ci sont définis par un Module inversible $\underline{L}$ \inserted{($\simeq g^*\mathcal{O}_{P'}(1)$)} sur P

\margin{102}
